\paragraph{A strict-reject, endpoint-accept repair from the stronger SPMV A0 state.}
The 0.2 ns current-STA window is now evaluated with a physical-delay objective and hard current-capacity constraints. The strict branch accepts no repair. The endpoint-Pareto branch accepts 6008/pin~12: physical delay decreases from 1.365955154541965 to 1.358316153998729 ns, two fixed setup/domain components improve and none regresses, while CPD remains exactly 7.034697713237392 ns. This is a legal plateau acceptance that strict CPD rejects. Resources change from zero to one SB pair, remain at 967 CB sources, and planar wirelength decreases from 148463 to 148445. The fixed initial 0.2 ns endpoint cohort still contains four components. No hold endpoint component changes.

After acceptance, the endpoint branch selects targets again from its own new STA. Its next legal joint candidate reduces 6997/pin~1 physical delay by approximately 0.185826 ns and improves that connection's native pin slack, but the entire fixed endpoint setup/hold vector and CPD remain unchanged. The endpoint rule therefore rejects this successor. The strict branch also rejects a larger joint candidate from its own unchanged initial state, despite two improved setup components, because CPD ties. Thus this completed operator demonstrates a secondary timing benefit and a resource tradeoff, but no later CPD gain unlocked by the plateau. Its shared two-policy flow takes 284.194 s, including three new native transactions; native transaction cores total 0.290 s. Candidate audit children total 11.114 s with the graph cache, while candidate FULL/STA wrappers still total 108.525 s. Matching earlier query results is explicit reuse, so this is not an independent cold-start comparison with the 199.891 s weighted-objective near-window flow.

\paragraph{Matched objective/capacity controls isolate adoption failure.}
All four cells use the identical unmodified A0 source, 6008/pin~12, retained tree scope, request bytes and current criticality. The completed weighted-soft and physical-hard cells are reused; only the missing cells run new native searches. Weighted cost with hard capacity still produces no physical-delay or endpoint improvement. Physical cost with soft congestion finds a path but fails current capacity feasibility; native occupancy is restored and no candidate is accepted. Physical cost with hard capacity produces the legal 7.639 ps connection gain and two improved setup components. This single-state factorial supports an interaction between adopting the physical objective and enforcing capacity during search. It does not establish a universal advantage across every connection or design. The two new cells take 32.486 s together; the weighted-hard candidate exactly matches a previously certified complete route, so its FULL/STA is reused with actual input/command/constraint identity checks. No failed soft-capacity transaction is presented as a legal timing result.

Actual adoption and downstream timing masking are distinct here. The weighted single repair did not shorten 6008/pin~12; the hard physical version did and changed endpoints. By contrast, the later 6997/pin~1 repair changes both actual routed delay and local native slack but no endpoint. The specific competing downstream paths have not yet been identified. The next registered diagnostic exports the existing native timing graph, edge delays and complete arrival/required tags at full precision to distinguish reconvergent path competition from mapping or constraint effects. It uses new necessary timing-cone observations of existing routes, not repeated routing to complete a report, and will derive any combined repair from that evidence.

\paragraph{Complete-route physical audits can also reuse exact graph evidence.}
A second metadata adapter preserves raw XML timing strings, directions, node/switch type counts and all fields expected by the unchanged complete-route auditor. It references the already bound edge/capacity cache; actual XML and every payload are checked on load. Its metadata build takes 29.172 s in addition to the earlier 33.122 s parent-cache build. The adapter explicitly supports the current graph's node kinds; CHANZ architectures are rejected rather than silently approximated. Candidate tree validation and native FULL/STA are not replaced by this cache.

In one contemporaneous validation, original XML extraction for the union of three existing routes takes 24.796 s, while cache load, digest checks and equivalent evidence construction take 3.310 s. Every raw graph-evidence field matches exactly. The unchanged complete-route auditor and physical-resource calculator then produce identical results on an accepted SPMV candidate, an accepted Tf candidate and a legal but timing-rejected R1 candidate, excluding audit wall time alone. Deliberately duplicated-net occupancy and an invalid switch receive identical rejection errors. No new routing or STA run is performed for these checks. These measurements support a data-access optimization, not an overall routing speedup. A separate cached FULL/STA wrapper has been prepared, preserving the native command and complete auditor, but has not yet been exercised on a fresh native candidate; its whole-process improvement remains unmeasured.
